\input header

\footline={\hss\tenrm D-\folio\hss}\pageno=1

{\bf Appendix D. Remarks Concerning Two Complex Variables.}\bigskip

Many of the arguments in these lectures are strictly
one-dimen\-sional. In fact several of our underlying principles
break down completely in the two variable case.\smallskip

In order to illustrate these differences, it is useful to consider the
family of (generalized)
{\bit H\'enon maps}, which can be described as follows. Choose a complex
constant
\[\delta\ne 0\] and a polynomial map \[f:\C\to\C\] of degree \[d\ge 2\],
and consider doubly infinite sequences of complex numbers \[\;\ldots\,,\,
z_{-1}\,,\,z_0\,,\,z_1\,,\,z_2\,,\,\ldots\;\] satisfying the recurrence
relation
$$	z_{n+1}-f(z_n)+\delta\,z_{n-1}\;=\;0\,.	$$
Evidently we can solve for \[(z_n\,,\,z_{n+1})\] as a polynomial function
\[\;F(z_{n-1}\,,\,z_n)\,\], where the transformation
\[\;\;F(z_{n-1}\,,\,z_n)\,=\,(z_n\,,\, f(z_n)-\delta z_{n-1})\;\;\] has Jacobian matrix
$$	\left[\matrix{0 & 1\cr -\delta & f'(z_n)\cr}\right] $$
with constant determinant \[\delta\], and with trace \[f'(z_n)\]. Similarly
we can solve for\break \[(z_{n-1}\,,\,z_n)\] as a polynomial function
\[F^{-1}(z_n\,,\,z_{n+1})\]. As an example, consider the quad\-ratic polynomial
\[f(z)=z^2+\lambda z\], having a
fixed point of multiplier \[\lambda\] at the origin. Then
$$	F(x\,,\,y)\;=\; (y\,,\; y^2+\lambda y-\delta x)\,,
	\eqno (*) $$
where \[x\] and \[y\] are complex variables. This map
has a fixed point at \[(0,0)\], where the eigenvalues \[\lambda_1\] and
\[\lambda_2\] of the Jacobian matrix satisfy \[\lambda_1+\lambda_2=\lambda\]
and \[\lambda_1\lambda_2=\delta\]. Evidently, by the appropriate choice
of \[\lambda\] and \[\delta\], we can realize any desired non-zero
\[\lambda_1\] and \[\lambda_2\]. If \[\lambda_1\ne\lambda_2\], then we
can diagonalize this Jacobian matrix by a linear change of coordinates.
\smallskip

{\QP{\bf Lemma D.1.} \it Consider any holomorphic transformation
\[F(x,y)=(x',y')\] in two complex variables, with
$$	x'=\lambda_1 x +{\cal O}
	(|x^2|+|y^2|)\,,\quad y'=\lambda_2 y+{\cal O}(|x^2|+|y^2|)	$$
as \[(x,y)\to(0,0)\].
If the eigenvalues \[\lambda_1\] and \[\lambda_2\] of the derivative
at the\break origin satisfy \[1>|\lambda_1|\ge
|\lambda_2|>|\lambda_1^2|\], then \[F\] is conjugate, under a local\break
holomorphic change of coordinates, to the linear map \[L(u,v)=(\lambda_1
u\,,\,\lambda_2v)\].\medskip}

{\bf Proof.} We must show that there exists a change of coordinates \[(u,v)=\phi(x,y)\],
defined and holomorphic throughout a neighborhood of the origin,
so that \[\phi\circ F\circ\phi^{-1}=L\]. As in the proof of the Ko$\!$enigs
Theorem, \S6.1, we first choose a constant \[c\] so that\break
\[1>c>|\lambda_1|\ge|\lambda_2|>c^2\]. To any orbit
$$	(x_0\,,\,y_0)\buildrel F\over\mapsto(x_1\,,\,y_1)
	\buildrel F\over\mapsto\cdots	$$
near the origin, we associate the sequence of points
$$	(u_n\,,\,v_n)\;=\;L^{-n}(x_n\,,\,y_n)\;=\;(x_n/\lambda_1^n\,,\,y_n
	/\lambda_2^n)	$$
and show, using Taylor's Theorem, that it converges geometrically to the
required limit \[\phi(x_0\,,\,y_0)\], with
successive differences bounded by a constant times \[(c^2/\lambda_2)^n\].
Details will be left to the reader.\QED

{\bf Remarks.} Some such restriction on the eigenvalues is essential.
As an example, for the map
$$	f(x,y)\;=\;(\lambda x,\,\lambda^2 y+x^2)\,,	$$
with eigenvalues \[\lambda\] and \[\lambda^2\],
there is no such local holomorphic change of coordinates. (Problem D-1.)
For a much sharper statement as to when linearization is possible,
even when one or both eigenvalues lie on the unit circle, compare Zehnder.


\smallskip

Now consider a H\'enon map \[F:\C^2\buildrel \approx\over\to\C^2\] as
in $(*)$ above,
with eigenvalues satisfying the conditions of D.1. Let \[\Omega\] be the
attractive basin of the origin. We claim that \[\phi\] extends
to a global diffeomorphism \[\Phi:\Omega\buildrel \approx\over\to\C^2\].
In fact, for any \[(x,y)\in\Omega\] we set \[\Phi(x,y)=L^{-n}\circ\phi\circ
F^{\circ n}(x,y)\]. If \[n\] is sufficiently large, then \[F^n(x,y)\] is
close to the origin, so that this expression is defined. Similarly,
\[\Phi^{-1}(u,v)=F^{-n}\circ\phi^{-1}\circ L^{\circ n}\] is well defined
for large \[n\]. This shows that
\[\Phi\] is a holomorphic diffeomorphism with holomorphic inverse.\smallskip

%***

%For suitable choice of \[\lambda_1\] and \[\lambda_2\] in the unit disk,
%it can shown ??? that \[F\] is locally conjugate to the linear map
%\[\;L:(w_1\,,\,w_2)\mapsto(\lambda_1w_1\,,\,\lambda_2w_2)\,\]. {\it Since
%\[F\] is a
%global diffeomorphism, it follows easily that the attractive basin \[B\]
%of the origin is analytically diffeomorphic to the entire coordinate space
%\[\C^2\] under a diffeomorphism \[\phi:B \buildrel \approx\over\longrightarrow
%\C^2\] which satisfies \[\phi\circ F=L\circ\phi\].} (Compare \S6.2.)
%\smallskip

%***

Note that this basin \[\Omega\] is not the entire space \[\C^2\]. For example, if \[
|z_1|\] is sufficiently large compared with \[|z_0|\], and if
$$	(z_0\,,\,z_1)\buildrel F\over\mapsto (z_1\,,\,z_2)\buildrel
	 F\over\mapsto (z_2\,,\,z_3)\buildrel F\over\mapsto\cdots	$$
then it is not difficult to check that \[\;|z_1|<|z_2|<|z_3|<\cdots\,\],
so that \[(z_0\,,\,z_1)\] is not in \[\Omega\].
{\it Thus we have constructed:\smallskip

$(1)$ a proper subset \[\Omega\subset\C^2\] which is analytically diffeomorphic to
all of \[\C^2\], and\smallskip

$(2)$ a non-linear map with an attractive basin which contains no critical
points.}\smallskip

\noindent Evidently neither phenomenon can occur in one complex variable.
Open sets satisfying (1) are called {\bit Fatou-Bieberbach domains} since they
were first constructed by Fatou, by an easy argument similar to that given
here, and then later independently by Bieberbach, who had a much more
difficult construction.\smallskip

The proof that there are only finitely many attracting cycles
also breaks down in two variables. Compare Newhouse.
\lln

{\bf Problem D-1.} For the map \[F(x,y)=(\lambda x\,,\,\lambda^2 y+x^2)\],
where \[\lambda\ne 0\,,\,1\], show that there is only one smooth
\[F$-invariant curve through the origin, namely \[x=0\]. By way of contrast,
for the associated linear map \[L(x,y)=(\lambda x\,,\,\lambda^2 y)\]
note that there are infinitely many \[F$-invariant curves \[y=cx^2\].
Conclude that \[F\] is not locally holomorphically conjugate to a linear map.

\vfil\eject % \bigskip\bigskip
\footline={\hss\tenrm E-\folio\hss}\pageno=1

\centerline{\bf Appendix E. Branched Coverings and Orbifolds.}\bigskip

This will be an outline of definitions and results, without proofs.
(See the list of references at the end.)
We will use ``branch point'' as a synonym for ``critical point''
and ``ramified point'' as a synonym for ``critical value''. Thus if
\[f(z_0)=w_0\] with \[f'(z_0)=0\], then \[z_0\] is called a {\bit branch
point} and the image \[f(z_0)=w_0\] is called a {\bit ramified point}.
More precisely, if
$$	f(z)\;=\;w_0+c(z-z_0)^n+({\rm higher\; terms})\,,	$$
with \[n\ge 1\] and
\[c\ne 0\], then the integer \[n=n(z_0)\] is called the {\bit branch index}
or the {\bit local degree} of \[f\] at
the point \[z_0\]. Thus \[n(z)\ge 2\] if \[z\] is a branch point,
and \[n(z)=1\] otherwise.

A holomorphic map \[p:S'\to S\] between Riemann surfaces is called a
{\bit covering map} if each point of \[S\] has a connected
neighborhood \[U\] which is {\bit evenly covered}, in that each connected
component of \[p^{-1}(U)\subset S'\] maps onto \[U\] by a conformal
isomorphism. A map \[p:S'\to S\] is {\bit proper} if the inverse image
\[p^{-1}(K)\] of any compact subset of \[S\] is a compact subset of \[S'\].
Note that every proper map is finite-to-one, and has a well
defined finite degree \[d\ge 1\]. Such a map may also be called a
\[d$-{\bit fold branched covering}. On the other hand, a covering
map may well be infinite-to-one. Combining these two concepts,
we obtain the following more general concept.\smallskip

{\bf Definition.} A holomorphic map \[p:S'\to S\] between Riemann surfaces
will be called a {\bit branched covering map} if every point of \[S\] has
a connected neighborhood \[U\] so that each connected component of \[p^{-1}
(U)\] maps onto \[U\] by a proper map.\smallskip

Such a branched covering is said to be {\bit regular} or
{\bit normal} if there exists a group \[\Gamma\] of conformal automorphisms
of \[S'\], so that two points \[z_1\] and \[z_2\] of \[S'\] have the same
image in \[S\] if and only if there is a group element \[\gamma\] with
\[\gamma(z_1)=z_2\]. In this case we can identify \[S\] with the quotient
manifold \[S'/\Gamma\]. In fact it is not difficult to check that the
conformal structure of such a quotient manifold is uniquely determined.
This \[\Gamma\] is called the group of {\bit deck transformations} of the
covering.

Regular branched covering maps have several special properties. For
example, each ramified point is isolated,
so that the set of all ramified points is a discrete subset of \[S\].
Furthermore, the branch index \[n(z)\] depends only on the target point
\[f(z)\], that is, \[n(z_1)=n(z_2)\]
whenever \[f(z_1)=f(z_2)\]. Thus we can define the {\bit ramification
function} \[\nu:S\to\{1,\,2,\,3,\,\ldots\}\] by setting \[\nu(w)\] equal
to the common value of \[n(z)\] for all points \[z\] in the pre-image
\[f^{-1}(w)\]. By definition, \[\nu(w)\ge 2\] if \[w\] is a ramified
point, and \[\nu(w)=1\] otherwise.\smallskip

{\bf Definition.} A pair \[(S,\nu)\] consisting of a Riemann surface \[S\]
and a ``ramification function'' \[\nu:S\to\{1,\,2,\,3,\,\ldots\}\] which
takes the value \[\nu(w)=1\] except at isolated points
will be called a Riemann surface {\bit orbifold}.

(Remark: Thurston's general concept of orbifold
involves a structure which is locally modeled on the quotient of a coordinate
space by a finite group. However, in the Riemann surface case only
cyclic groups can occur, so a simpler definition can be used.)

{\bf Definition.}  If \[S'\] is simply connected, then a regular branched
covering \[p:S'\to S\] with ramification function \[\nu\]
will be called the {\bit universal covering}
for the orbifold \[(S,\,\nu)\]. We will use the notation \[\tilde S_\nu\to S\]
for this universal branched covering. The associated
group \[\Gamma\] of deck transformations is called the {\bit fundamental
group} \[\pi_1(S,\,\nu)\] of the orbifold.

{\QP{\bf Lemma E.1.} \it With the following exceptions, every
Riemann surface orbifold
\[(S,\,\nu)\] has a universal covering surface \[\tilde S_\nu\]
which is unique up to conformal isomorphism over \[S\]. The only exceptions
are given by:

(1) a surface \[S\approx\hat\C\]
with just one ramified point, or

(2) a surface \[S\approx\hat\C\] with two ramified points
for which \[\nu(w_1)\ne\nu(w_2)\].

\noindent
In these exceptional cases, no such universal covering exists.\medskip}

By definition,
the {\bit Euler characteristic} of an orbifold \[(S,\,\nu)\] is
the rational number
$$      \chi(S,\,\nu)\;=\;\chi(S)+\sum \big({1\over\nu(w_j)}-1\big)\,, $$
to be summed over all ramified points, where \[\chi(S)\] is the usual
Euler characteristic of \[S\]. Intuitively speaking,
each ramified point \[w_j\] makes a contribution of \[+1\] to the usual
Euler characteristic \[\chi(S)\], but a smaller contribution of \[1/\nu(w_j)\]
to the orbifold Euler characteristic. Thus \[\chi(S,\,\nu)<\chi(S)\le 2\],
or more precisely
$$	\chi(S)-r\;<\; \chi(S,\,\nu)\;\le\;\chi(S)-r/2	$$
where \[r\] is the number of ramified points. {\it As an example, if
\[\chi(S,\,\nu)
\ge 0\], with at least one ramified point, then it follows that \[\chi(S)>0\],
so the base surface \[S\] can only be \[D,\,\C\] or \[\hat \C\], up
to isomorphism.} Compare E.5 below.

If there are infinitely
many ramified points, note that we must set \[\chi(S,\,\nu)=
-\infty\]. Similarly, if \[S\] is
a surface which is not of finite type, then we must set\break
\[\chi(S,\,\nu)=\chi(S)=-\infty\].\smallskip

If \[S'\] and \[S\] are provided with ramification functions \[\mu\] and
\[\nu\] respectively, then a branched covering map \[f:S'\to S\] is said
to yield
a {\bit covering map \[(S'\,,\,\mu)\to(S\,,\,\nu)\] between orbifolds}
if the identity
$$      n(z)\mu(z)\;=\;\nu(f(z))       $$ 
is satisfied for all \[z\in S'\],
% if \[n(z)\nu'(z)\] is a multiple of \[\nu(f(z))\] for each \[z\in S'\],
where \[n(z)\] is the branch index.
% Similarly, we say that \[f^{-1}\]
%is {\bit locally holomorphic} between orbifolds if \[\nu(f(z)\] is a
% multiple of \[n(z)\nu'(z)\], and that \[f\] is a {\bit covering map} of
As an example, the universal covering map \[\tilde S_\nu\to
(S\,,\,\nu)\] is always a covering map of orbifolds, where \[\tilde S_\nu\]
is provided with the trivial ramification function \[\mu\equiv 1\].

{\QP\bf Lemma E.2. {\it \[f:(S'\,,\,\mu)\to(S\,,\,\nu)\] is a covering map
between orbifolds if and only if it lifts to a conformal isomorphism from the
universal covering \[\tilde S'_\mu\] onto \[\tilde S_\nu\].
If \[f\] is a covering in this sense, and has finite degree \[d\],
then the Riemann-Hurwitz formula of \S5.1 takes the form
$$	\chi(S'\,,\,\mu)\;=\; \chi(S\,,\,\nu)d\,.		$$
In particular, if the universal covering of \[(S\,,\,\nu)\] is a covering
of finite degree \[d\], then \[\;\chi(\tilde S_\nu)=\chi(S,\,\nu)d\].}
\medskip}

\noindent The fundamental group and the Euler characteristic
% of any Riemann surface orbifold 
are related to each other as follows.

\smallskip

{\QP{\bf Lemma E.3.} \it Let \[(S,\,\nu)\] be any Riemann surface orbifold
which possesses a universal covering:

If \[\chi(S,\,\nu)>0\], then the fundamental
group \[\pi_1(S,\,\nu)\] is finite.

If \[\chi(S,\,\nu)=0\], then the fundamental group contains either \[\Z\]
or \[\Z\oplus\Z\] as a subgroup of finite index.

If \[\chi(S,\,\nu)<0\], then the fundamental
group contains a non-abelian free product \[\Z *\Z\], and hence cannot contain
any abelian subgroup of finite index.\medskip}

\noindent
The Euler characteristic and the geometry of \[\tilde S_\nu\] are related
as follows.

{\QP{\bf Lemma E.4.} \it Suppose that \[S\] is obtained from a compact
Riemann surface by
removing at most a finite number of points. Then \[\tilde S_\nu\] is
either spherical, hyperbolic or Euclidean according as \[\chi(S\,,\,\nu)\]
is positive, negative or zero.\medskip}

{\bf Remark.} In all other cases, that is whenever \[S\] is
not isomorphic to a finitely punctured compact surface, \[S\] is
hyperbolic, and it follows that the universal covering
\[\tilde S_\nu\] must also be hyperbolic.

{\bf Examples.} If \[S=\C\] with two ramified
points \[\nu(1)=\nu(-1)=2\], then the map\break
\[z\mapsto\cos(2\pi z)\] provides
a universal covering \[\C\to (\C,\,\nu)\]. The Euler characteristic
\[\chi(\C,\,\nu)\] is zero, and the fundamental group \[\pi_1(S,\,\nu)\]
consists of all transformations of the form \[\gamma:z\mapsto  n\pm z\].

For \[S=\hat\C\] with three ramified points \[\nu(0)=\nu(1)=\nu(
\infty)=2\],
the rational map \[\pi(z)=-4z^2/(z^2-1)^2\] provides a
universal covering \[\hat\C\to(\hat\C,\,\nu)\]. In this case we have
\[\chi(\hat\C,\,\nu)=1/2\], and
the degree is equal to \[\chi(\hat\C)/\chi(\hat\C,\,\nu)=4\].
The fundamental group consists of all transformations \[\gamma:z\mapsto
\pm z^{\pm 1}\].

For \[S=\hat\C\] with four ramified points of index \[\nu(w_j)=2\], then the
torus \[T\] described in \S5 provides a regular 2-fold branched covering.
Its universal covering \[\tilde T\] can be identified with the
universal covering of \[(\hat\C,\,\nu)\]. In this case, the Euler
characteristic \[\chi(\hat\C,\,\nu)\] is zero.\smallskip


{\bf Remark E.5.} There are relatively few cases in which \[\chi(S\,,\,
\nu)\ge 0\]. In fact all of these cases can be listed quite explicitly
as follows. The unramified cases are very well known, namely
the sphere, plane or disk with \[\chi>0\]; and the punctured
plane or disk, and the infinite families consisting of annuli and tori,
all with \[\chi=0\].

 By the ``ramification indices'' we will
mean the list of values of the ramification function at the \[r\]
ramified points,
ordered for example so that \[\nu(w_1)\le\cdots\le\nu(w_r)\].

If \[\chi(\hat\C,\,\nu)>0\] with \[r>0\],
then the ramification indices must be either
\[(n,n)\] or \[(2,2,n)\] for some \[n\ge 2\], or \[(2,3,3)\,,\, (2,3,4)\]
or \[(2,3,5)\]. These five possibilities correspond to the five types of
finite rotation groups of the 2-sphere; namely to the cyclic,
dihedral, tetrahedral, octahedral and icosahedral groups respectively.

If \[\chi(\hat\C,\,\nu)=0\], then the ramification indices must be either
\[(2,4,4)\], \[(2,3,6)\], \[(3,3,3)\] or \[(2,2,2,2)\].
These correspond to the automorphism groups of the tilings of \[\C\] by
squares, equilateral triangles, alternately colored equilateral triangles,
and parallelograms respectively. In the parallelogram case, note that there
is actually a one
complex parameter family of distinct possible shapes, corresponding to the
cross-ratio of the four ramified points.

Similarly, if \[\chi(\C,\,\nu)\] or \[\chi(D,\,\nu)\] is strictly positive,
then we must have \[r\le1\]; while if \[\chi(\C,\,\nu)\] or \[\chi(D,\,\nu)\]
is zero, then we must have \[r=2\] with ramification  indices \[(2,2)\].
This is the complete list.


\vfill
\eject % \bigskip\bigskip


\footline={\hss\tenrm F-\folio\hss}\pageno=1

\centerline{\bf Appendix F. Parameter Space.}\bigskip

A very important part of complex dynamics, which has barely been mentioned in
these notes, is the study of parametrized families of mappings. As an example,
consider the family of all quadratic polynomial maps. A priori, a quadratic
polynomial is specified by three complex parameters; however
any such polynomial can be put into the unique normal form
$$	f(z)\;=\;z^2+c \eqno (1)	$$
by an affine change of coordinates. (Other normal forms which have been used
are\break \[	\omega\;\mapsto\;\omega^2+\lambda \omega\],
with a preferred fixed point of multiplier \[\lambda\] at the origin, or
$$	w\; \mapsto\;\lambda w (1-w) \eqno (2) $$
which is more or less equivalent provided that
\[\lambda\ne 0\]. Here \[4c=\lambda(2-\lambda)\] and\break
 \[\omega=z-\lambda/2=-\lambda w\].) Using such a normal form, we can make
a computer picture in the
{\bit parameter space} consisting of all complex constants \[c\] or \[\lambda
\]. Each pixel in such a picture, corresponding to a small square in the
parameter space, is to be assigned some color, perhaps only black or white,
which depends on the dynamics of the corresponding quadratic map.

The first crude pictures of this type were made by Brooks and Matelski,
as part of a study of Kleinian groups. They used the normal form (1), and
introduced the open set consisting
of all points of the \[c$-plane for which the corresponding quadratic map
has an attracting periodic orbit in the finite plane.
I will use the notation \[{\cal H}\] (standing
for ``hyperbolic'') for this Brooks-Matelski set.
At about the same time,
Hubbard (unpublished) made much better pictures of a quite different
parameter space arising from Newton's method for cubics. Mandelbrot,
perhaps inspired by Hubbard, made corresponding pictures for quadratic
polynomials, using the normal form (2) and also a variant of (1).
In order to avoid confusion, let me translate all of Mandelbrot's
definitions to the normal form (1).
He introduced two different sets, which I will call \[Q\subset
M\]. Mandelbrot did not give these sets
different names, since he believed that they were identical. By definition,
a parameter value \[c\] belongs to \[Q\] if the corresponding filled Julia set
contains an interior point, and belongs to \[M\] if its filled Julia
set contains the critical point \[z=0\] (or equivalently is connected).
The Brooks-Matelski set
\[{\cal H}\] is a subset of \[Q\subset M\]. Mandelbrot made quite good computer
pictures, which seemed to show a number of isolated ``islands''. Therefore,
he conjectured that \[Q\] [or \[M\]] has many distinct connected components.
(The editors of the journal thought that his islands were specks of dirt,
and carefully removed them from the pictures.) Mandelbrot also described
an important smaller set \[{\cal L}\subset Q\] which he believed to be the
principal connected component of \[Q\]. This set \[{\cal L}\]
consists of a central cardioid
\[{\cal L}_0\] with some (but not all) boundary points included,
together with countably many smaller nearly-round disks which are pasted on
inductively, in an explicitly described pattern.\smallskip

Although Mandelbrot's statements in this first paper
were not completely right, he deserves a great
deal of credit for being the first to point out the extremely complicated
geometry associated with the parameter space for quadratic maps.
His major achievement has been to demonstrate
to a very wide audience that such complicated ``fractal''
objects play an important role in a number of mathematical sciences.
\smallskip

The first real mathematical breakthrough came with Douady and Hubbard's work in
1982. They introduced the name {\bit Mandelbrot set} for the compact set
\[M\] described above, and
provided a firm foundation for its mathematical study,
proving for example that \[M\] is connected with connected complement.
(Meanwhile,
Mandelbrot had decided empirically that his isolated islands were actually
connected to the mainland by very thin filaments.)
Already in this first paper, they showed that each hyperbolic component
of the interior of \[M\] can be canonically
parametrized, and showed that the boundary \[\partial M\] can be profitably
studied by following external rays.\smallskip

It may be of interest to compare the three sets \[{\cal H}\subset Q\subset M\]
in parameter space. They are certainly different since \[{\cal H}\] is open,
\[M\] is compact, and \[Q\] is neither.
Using Sullivan's work (\S13), we can say that
\[Q\] consists of \[{\cal H}\] together with a very sparse set of boundary
points, namely those for which the corresponding map has
either a parabolic orbit or
a Siegel disk. Quite likely, there is no difference between these three sets
as far as computer graphics are concerned, since it is widely
conjectured{\bf *} that the Brooks-Matelski set \[{\cal H}\]
is equal to the interior of \[M\], and that
\[M\] is equal to the closure
of \[{\cal H}\]. However, as far as practical computing is concerned, it should
be noted that it is quite easy to
test (at least roughly)
whether a parameter value belongs to \[M\], but somewhat harder to decide
whether it belongs to \[{\cal H}\] (compare \S10.5), and very difficult
to decide
whether it belongs to \[Q\]. (Compare Appendix G. Here I am only speaking of
approximate tests. To decide precisely whether a given point belongs to
\[M\] or to whether a given point belongs to \[{\cal H}\] may be very
difficult. As a specific example, I have no idea whether or not the point
\[c=-1.5\] belongs to \[{\cal H}\].)\smallskip

Another important development came in 1983, with
the work of Ma\~n\'e, Sad and Sullivan on stability of the Julia set
\[J(f)\] under deformation of \[f\]. (Compare the discussion following \S14.3.)
These results were obtained independently by Lyubich.
The study of parameter space for
higher degree polynomials began some five years later with the work of
Branner and Hubbard. Using the normal form
$$	f(z)\;=\;z^3-3a^2z+b\,, $$
with the two critical points at \[z=\pm a\], they proved that the {\bit cubic
connectedness locus}, consisting of all parameter pairs \[(a,b)\] for which
\[J(f)\] is connected, is a {\bit cellular set}. (Compare \S17.3.)
In particular, this set is compact and connected. A corresponding result for
polynomials of higher degree has recently been obtained by Lavaurs.
Parameter space studies for rational maps are more awkward,
since there is no obvious normal form. However, an important beginning
has been made by Rees. All of these studies are very new, and much remains
to be done.\smallskip

Here are two problems for the reader.
\vfil
\noindent ------------

\noindent {\bf *} Douady and Hubbard have shown that these conjectures are true
if the set \[M\] is locally connected. Recent work of Yoccoz lends very
strong support to the belief that \[M\] is indeed locally connected.
\eject

{\bf Problem F-1.} Show that every polynomial map of degree \[d\ge 2\]
is conjugate, under an
affine change of coordinates, to one in the ``Fatou normal form''
$$	f(z)\;=\; z^d + a_{d-2}z^{d-2}+\cdots+a_1z+a_0\,.	$$
Let \[P(d)\cong\C^{d-1}\] be the space of all such maps.
Show that the cyclic group \[Z(d-1)\] of \[(d-1)$-st roots of unity acts on
\[P(d)\] by linear conjugation, replacing \[f(z)\] by \[f(\omega z)/\omega\],
and show that the quotient \[P(d)/Z(d-1)\] can be identified with the
``moduli space'' of degree \[d\] polynomials up to affine conjugation.
If \[d\ge 4\], show that this moduli space is not a manifold.\smallskip

{\bf Problem F-2.} Show that every quadratic rational map is conjugate, under
a fractional linear change of coordinates, for example to one in the form
$$	f(z)\;=\;1+(z^2-1)/(az^2-b) $$
where \[a\ne b\],
with critical points at zero and \[\infty\] and a fixed point at \[z=1\].
(In general this form is not unique, since such a
map actually has three fixed points, and since the roles of zero and infinity
can be interchanged.)% Show that the moduli space of quadratic rational maps
%up to fractional linear change of coordinates has just one non-manifold point,
%namely \[z\mapsto 1/z^2\].\smallskip


\vfil\eject
\footline={\hss\tenrm G-\folio\hss}\pageno=1

\centerline{\bf Appendix G. Remarks on Computer Graphics.}\medskip

In order to make a computer picture of some complicated compact set
\[L\subset\C\], for example a Julia set,
we must compute a matrix of small
integers, where the \[(i,j)$-th entry describes the color (perhaps
only black or white) which is assigned to the \[(i,j)$-th ``pixel'' of
the computer screen. Each pixel represents a small square in the complex plane,
and the color which is assigned must tell
us something about intersection of \[L\] with this square.

In the case of a quadratic Julia set \[J(f)\], one very fast method
involves following iterates of the inverse map \[f^{-1}\],
taking all possible branches. (Compare \S3.9.) As Mandelbrot points out,
this method yields an excellent
picture of the outer parts of the \[J(f)\] but shows very little detail in
the inner parts. If we think of the electrostatic field produced by
an electric charge on \[J(f)\], this method will emphasize only those parts of
the Julia set at which ``lines of force'' (or ``external rays'' in the language
of \S18) tend to land.

A slower but much better procedure for plotting \[J(f)\] involves
iterating the map \[f\] for some large number (perhaps 50 to 50000) of times,
starting at the midpoint of each pixel. If the orbit ``escapes'' from
a large disk after \[n\] iterations, then the corresponding pixel is
assigned a color which depends on \[n\]. In more refined versions of
this method, one computes not only the value of the \[n$-th iterate of \[f\]
but also the absolute value of its derivative. Compare the discussion
below. Similar remarks apply to the {\bit Mandelbrot set} \[M\], as defined
by Douady and Hubbard. (Compare
Appendix F). In this case one takes the quadratic map
corresponding to the midpoint of the square and follows the orbit of its
critical point.

{\bf Remark.}
In order to understand some of the limitations of this method, consider
the situation near a fixed point \[f(z_0)=z_0\] in the Julia set.
First consider the repelling
case, with multiplier say of absolute value two. If we start at at point
\[z\] at distance \[1/1000\] from \[z_0\], then the distance from
\[z_0\] will roughly double with each iteration. Hence, after only ten
iterations the image of \[z\] will move substantially away from \[z_0\].
The result will be a computer picture which is quite sharp and accurate
near \[z_0\]. (Figures 3, 4.)

Now suppose that we try to construct a picture for \[z\mapsto z+z^4\] by
the same method. Again start with a point \[z\in J\] at a distance of
\[\epsilon\approx 1/1000\] from the fixed point at zero. Examining
the proof of 7.2, we see that the associated coordinate \[w_0=- 1/
3z_0^{\,3}\], which increases by one under each iteration, is roughly
equal to \[-1/3\epsilon^3\approx 300,000,000\]. {\it Thus
we would have to follow such an orbit for some \[300,000,000\]
iterations in order
to escape from a neighborhood of \[z=0\].} The result in practice will be a
false picture which shows everything near the origin to be in the Fatou set.
(This difficulty was eliminated in Figure 8 by using a special computer
program, which extrapolated iterates of \[f\] in order to distinguish
different Fatou components. Similarly, Figures
10, 12 were plotted with special purpose programs.)

For the fixed points of Cremer type, the situation is much worse. As far
as I know, no useful
computer picture of such a point has ever been produced!\smallskip

Many Julia sets are made up of very fine filaments. For such sets, it
is essential to make some kind of distance estimate in order to obtain
a sharp picture. For all of the center points of our pixels will quite likely
lie outside of these filaments, and hence correspond to escaping orbits.
But a good distance estimate can tell us that our pixel intersects the
\break

\noindent
set \[J(f)\], even though its center point is outside. In the case of
the Mandelbrot set \[M\], this procedure is even more important, since \[M\]
contains both large regions and also very fine filaments. Indeed,
it was precisely the difficulty of seeing such filaments which led
to Mandelbrot's incorrect belief that \[M\] has many components.

Here is an example of how first derivatives can be used to make such
distance estimates, following Fisher.
Consider a rational map \[f:\hat\C\to\hat\C\] with a superattractive
fixed point at the origin. Let \[\Omega\] be the basin of attraction
of this fixed point. Let us assume, to simplify the discussion, that this
basin is connected, simply connected,\break and
contains no other critical point of \[f\]. Then it is not difficult to
show that the B\"otk$\!$her coordinate of \S6.7 can be defined throughout
\[\Omega\], and yields a conformal isomorphism
\[\;\phi:\Omega\to D\;\] with \[\;\phi(f(z))=\phi(z)^n\,\].
Define the {\bit canonical potential function}\break
\[G:\Omega-\{0\}\to \R\;\] by
$$	G(z)=\log|\phi(z)|<0\,.	$$
(Compare \S17.) Denoting the gradient vector of \[G\] by \[G'\], we will show
that:\smallskip

(1) {\it the function \[G\] and the norm \[\| G' \|= |\phi'(z)/
\phi(z)|\] are easy to compute, and}\smallskip

(2) {\it the distance of \[z\] from the boundary of \[\Omega\]
can be computed, up to a factor\break
\indent of two, from a knowledge of \[G\] and \[\| G' \|\].}\smallskip

\noindent
In fact, for any orbit \[z_0\mapsto z_1\mapsto\cdots\] in \[\Omega\], it is easy
to check that
$$	G(z_0)=\lim_{k\to\infty}\log|z_k|/n^k\,.	$$
Since the convergence is locally uniform, we can also write
$$	\|G'(z_0)\|=\lim_{k\to\infty}{|dz_k/dz_0|\over n^k|z_k|}\,.	$$
In both cases, the successive terms can easily be computed
inductively, and
we obtain good approximations by iterating until \[|z_k|\]
is small. (If many iterations do not yield any small \[z_k\], then we must
assume that \[z_0\ne \Omega\], and set \[G=0\].)

Setting \[\phi(z)=w\], a brief computation shows that
the \Poincare metric on \[\Omega\] can be written as
$$	{2|dw|\over 1-|w|^2}\;=\; {2|\phi'(z)dz|\over 1-|\phi(z)|^2}\;
	=\;{\|G'(z)dz\|\over |\sinh G(z)|}\;.	$$
As an immediate consequence of the Quarter Theorem, \S A.7, A.8, we
obtain the following.

{\QP {\bf Corollary.} {\it The distance between \[z\] and the
boundary of \[\Omega\] is equal to\break
\[|\sinh( G)|/\|G'\|\], up to a factor of two.}\medskip}

\noindent If \[z\] is very close to \[\partial \Omega\], then \[G\] is small and
this distance estimate is very close to the ratio \[|G|/\|G'\|\]. It is
interesting to note that this is just the step size which would be
prescribed if we tried to solve the equation \[G(z)=0\] by Newton's method.
\smallskip

There are similar estimates in the case of a superattractive fixed point
at infinity, or in other words for the basin \[\C\ssm K(f)\] of a polynomial
map. For further information, the reader is referred to
Fisher, to Milnor [M5], and to Peitgen.

\vfill
\rightline{jwm, version of 9-4-91}
\end



